\section*{Chapitre \ref{ch:ml:trees-and-boosting} --- Arbres et boosting}

\begin{solution}{exo:ml:trees-and-boosting:1}
$0.3\sigma^2 + 0.7\sigma^2/50 = 0.314\sigma^2$~; avec une infinité d'arbres, $0.3\sigma^2$~: seule la
décorrélation abaisse le plancher.
\end{solution}

\begin{solution}{exo:ml:trees-and-boosting:2}
$400\times(-0.1 - 0.08)^2 + 600\times(0.2 - 0.08)^2 = 12.96 + 8.64 = 21.6$.
\end{solution}

\begin{solution}{exo:ml:trees-and-boosting:3}
$F_1 = 0.1\times0.8 = 0.08$, résidu 0.92. Si chaque arbre prédisait exactement le résidu courant, chaque étape conserverait
une fraction $1 - \eta = 0.9$ de celui-ci~: $0.9^{10} = 0.349$ après dix étapes.
\end{solution}

\begin{solution}{exo:ml:trees-and-boosting:4}
Une ligne hors sac partage son mois, donc ses rendements des facteurs, avec des centaines de lignes dans le sac aux caractéristiques proches~;
les arbres ont appris de ces lignes les mouvements du facteur de ce mois, si bien que la ligne n'est pas hors échantillon. Les mois de test ne partagent
rien avec les mois d'entraînement.
\end{solution}

\begin{solution}{exo:ml:trees-and-boosting:5}
Les valeurs par défaut donnent en moyenne 0.09~\% avec un écart type de 0.07 points selon les graines~; 0.14~\% sur une seule graine se situe à moins
d'un écart type de leur moyenne. Réévaluer le candidat avec les mêmes quatre graines (et sur des blocs purgés) avant de le déclarer
meilleur.
\end{solution}

\begin{solution}{exo:ml:trees-and-boosting:6}
Le nombre d'arbres a été choisi sur les mois de test, si bien que le score de test est un score de sélection, optimiste de
la \cref{prop:ml:validation:max}. Il faut arrêter l'apprentissage sur un bloc de validation purgé à la fois des mois d'entraînement et des mois de test.
\end{solution}

\begin{solution}{exo:ml:trees-and-boosting:7}
\texttt{monotone(ofi5\_sign=-1)}~: R-carré de 0.305, contre 0.333 sans contrainte et 0.364 avec le bon signe~; augmenter
le déséquilibre fait désormais baisser la prévision pour 95~\% des lignes. Une contrainte fausse est un a priori faux imposé exactement~; elle coûte
plus qu'aucune contrainte.
\end{solution}

\begin{solution}{exo:ml:trees-and-boosting:8}
Avec une perte quadratique et des ajustements exacts, $F_m = F_{m-1} + \eta(y - F_{m-1})$, donc $y - F_m = (1 - \eta)(y - F_{m-1})$ et, à partir
de $F_0 = 0$, $y - F_M = (1 - \eta)^My$. La progression de l'ajustement dépend à peu près de $\eta M$ (pour $\eta$ petit, $(1 - \eta)^M
\approx e^{-\eta M}$)~: diviser par deux le \omterm{def:ml:trees-and-boosting:boosting}{taux d'apprentissage} double le nombre d'arbres nécessaires, et arrêter tôt à un $\eta$ donné revient à
choisir un $\eta M$ effectif plus petit. Les arbres réels ne suivent les résidus qu'en partie et ajustent aussi le bruit, ce qui explique qu'un petit
$\eta$ avec davantage d'arbres généralise mieux qu'un grand avec peu d'arbres.
\end{solution}

\begin{solution}{pb:ml:trees-and-boosting:1}
\begin{enumerate}
\item 0.24~\%, $-0.20~\%$, $-3.42~\%$~: des arbres plus profonds ajustent davantage de bruit (le terme de variance).
\item 0.20~\% hors échantillon, 0.42~\% hors sac, 2.76~\% dans l'échantillon.
\item Les lignes hors sac partagent les rendements des facteurs de leur mois avec les lignes dans le sac.
\item Davantage d'arbres supprime le terme $(1-\rho)\sigma^2/B$~; seuls des arbres moins corrélés abaissent $\rho\sigma^2$.
\item $-0.19~\%$.
\item Après 34 arbres à 0.1 (test 0.05~\%) et 233 arbres à 0.01 (test 0.26~\%).
\item À 0.1, le R-carré de test tombe à $-3.38~\%$~; à 0.01, à 0.03~\%.
\item Chaque arbre déplace peu l'ajustement, si bien que le bruit doit être ajusté lentement et que la courbe de validation est plate près de son maximum.
\item De $-0.54~\%$ à 0.39~\%.
\item Quatre feuilles, \omterm{def:ml:trees-and-boosting:boosting}{taux d'apprentissage} 0.02, 200 arbres, 20 lignes par feuille~; 0.36--0.39~\% avec quatre graines, 0.32~\% sur le test.
\item Écarts types de 0.07 points pour les valeurs par défaut, 0.014 pour le candidat retenu~: cinq fois plus bruyant.
\item Quatre feuilles, 1\,000 lignes par feuille, sans sous-échantillonnage~: 0.35~\% en validation, 0.31~\% en test.
\item R-carré de 0.333 à 0.364~; violations passant de 7.0~\% des lignes à aucune.
\item Elle supprime les directions d'ajustement que seul le bruit soutient~: une restriction que la vérité satisfait abaisse la variance
sans ajouter de biais.
\item \emph{Résultat nommé}~: le score de validation des valeurs par défaut varie avec la graine de 0.07 points (0.00--0.17~\%) et celui du
candidat retenu de 0.014, contre un gain du réglage d'environ 0.3 points (de 0.09~\% à 0.39~\%)~; le candidat retenu perd 0.07 points
de la validation au test.
\item Le choix du plateau~: aussi bon sur le test, moins flexible et moins sensible à la graine.
\item Trois à cinq, assez pour voir un écart type~; davantage si les candidats sont à moins d'un écart type les uns des autres.
\item Chaque configuration et son score, les graines, le bloc de validation et la règle utilisée pour choisir.
\item L'échantillon effectif pour les effets des caractéristiques est celui des mois, si bien qu'ajouter des actions aide moins qu'ajouter des mois~;
les configurations complexes resteraient bruitées.
\item Des \omterm{def:ml:trees-and-boosting:boosting}{taux d'apprentissage} faibles, peu de feuilles, beaucoup de lignes par feuille, un \omterm{def:ml:trees-and-boosting:early}{arrêt précoce} sur des données purgées, plusieurs graines et un choix
pris sur le plateau.
\end{enumerate}
\end{solution}

\begin{solution}{iq:ml:trees-and-boosting:1}
Le \omterm{def:ml:trees-and-boosting:bagging}{bagging} moyenne des modèles ajustés indépendamment sur des échantillons bootstrap, ce qui réduit la variance~; le boosting ajuste des modèles en séquence,
chacun sur les erreurs de l'ensemble constitué jusque-là, ce qui réduit le biais (et, avec le shrinkage et l'\omterm{def:ml:trees-and-boosting:early}{arrêt précoce}, contrôle la variance).
\iqlookfor{parallèle contre séquentiel, variance contre biais.}
\end{solution}

\begin{solution}{iq:ml:trees-and-boosting:2}
Ceux qui contrôlent la variance~: un \omterm{def:ml:trees-and-boosting:boosting}{taux d'apprentissage} faible avec \omterm{def:ml:trees-and-boosting:early}{arrêt précoce}, peu de feuilles ou une faible profondeur, un minimum élevé de
lignes par feuille, sous-échantillonnage des lignes et des colonnes, une pénalité $\ell_2$. Tous vont dans le sens d'une moindre flexibilité.
\iqlookfor{les régularisations d'abord, et le sens à donner sur des données à faible rapport signal/bruit.}
\end{solution}

\begin{solution}{iq:ml:trees-and-boosting:3}
Les lignes hors sac partagent les dates (ainsi que les chocs communs de facteurs et les étiquettes qui se chevauchent) avec les lignes dans le sac, si bien qu'elles ne sont pas des
données de test indépendantes~; le score est optimiste. Utiliser une validation purgée et ordonnée dans le temps.
\iqlookfor{la dépendance entre les lignes de données de panel.}
\end{solution}

\begin{solution}{iq:ml:trees-and-boosting:4}
Une contrainte imposant que le modèle soit non décroissant (ou non croissant) en une variable. Candidats~: le déséquilibre du flux d'ordres (OFI) et
le déséquilibre des files (une pression plus forte côté acheteur ne doit pas prédire un prix plus bas)~; des frais d'emprunt dans un modèle de short squeeze~; une
distance à une limite de variation de cours.
\iqlookfor{la définition et des signes justifiés économiquement.}
\end{solution}

\begin{solution}{iq:ml:trees-and-boosting:5}
Comparer d'abord au bruit~: réévaluer les deux avec plusieurs graines et sur des blocs purgés~; s'ils sont à moins d'un écart
type, retenir le plus simple, celui qui génère moins de rotation ou le mieux contraint.
\iqlookfor{le bruit avant le classement, puis la simplicité.}
\end{solution}

\begin{solution}{iq:ml:trees-and-boosting:6}
Le gradient négatif de $\frac12(y - F)^2$ par rapport à $F$ vaut $y - F$, soit le résidu, si bien que chaque arbre ajuste des résidus. Avec la perte de
Huber, le gradient est le résidu tronqué à $\pm\delta$~: les grandes erreurs (rendements à queues épaisses) tirent moins l'ajustement.
\iqlookfor{le calcul du gradient et la robustesse de la perte de Huber.}
\end{solution}
